% META: {"id":"TCOMPL-LOG-CR-012","chapitre":"TCOMPL-LOGARITHME-HISTORIQUE","type_objet":"cours","capacites_codes":["C5"],"sources_inspiration":[],"mode_creation":"ex_nihilo","status":"approved","origin":"generated","status_history":["generated","audited","accepted","approved"],"release_acceptance":"RELEASE_OWNER_BATCH_ACCEPTANCE","acceptance_closure_digest":"sha256:9b3ccf9a81c5520fbb7e03b7b2d3e7bf2057a02834b6d908bf3be5f49f3b553f","acceptance_receipt_digest":"sha256:4fad86f0282e0fa4e0419cca1ad6379ae7af5a280c04a4a90880f90a1c044242"}

\section{Dérivée du logarithme népérien}

% ============================================================
% STRATE 1 — Démonstration de la dérivée
% ============================================================

On admet que la fonction $\ln$, réciproque de la fonction exponentielle, est dérivable sur $]0,+\infty[$.

\theoreme[Dérivée de $\ln$]{
Pour tout $x > 0$ :
\[
  \ln'(x) = \frac{1}{x}.
\]
}

\demonstration{
Pour tout $x>0$, on a $\mathrm{e}^{\ln x} = x$. En dérivant les deux membres de cette égalité (le membre de gauche par la formule de dérivation d'une composée, en admettant la dérivabilité de $\ln$) :
\[
  \ln'(x) \times \mathrm{e}^{\ln x} = 1,
\]
soit $\ln'(x) \times x = 1$ (car $\mathrm{e}^{\ln x}=x$), d'où $\ln'(x) = \dfrac{1}{x}$.
}

% BEGIN-VERIFY
% from sympy import symbols, ln, diff
% x = symbols('x', positive=True)
% assert diff(ln(x), x) == 1/x
% END-VERIFY

% ============================================================
% STRATE 2 — Dérivée de fonctions composees avec ln
% ============================================================

\propriete[Dérivée de $\ln u$]{
Si $u$ est une fonction dérivable et strictement positive sur un intervalle $I$, alors $\ln u$ est dérivable sur $I$ et :
\[
  (\ln u)'(x) = \frac{u'(x)}{u(x)}.
\]
}

\exemple{
Pour $f(x) = \ln(x^2+1)$, on a $u(x)=x^2+1$ (toujours strictement positive), $u'(x)=2x$, donc :
\[
  f'(x) = \frac{2x}{x^2+1}.
\]
}

% BEGIN-VERIFY
% from sympy import symbols, ln, diff, simplify
% x = symbols('x')
% f = ln(x**2 + 1)
% fp = diff(f, x)
% assert simplify(fp - 2*x/(x**2+1)) == 0
% END-VERIFY

\erreurFrequente{
\textbf{Confondre $(\ln u)'$ et $\ln(u')$.} La dérivée de $\ln u$ est $\dfrac{u'}{u}$ (un quotient), pas le logarithme de la dérivée de $u$. Cette dernière expression n'a d'ailleurs de sens que si $u'$ est elle-même strictement positive, ce qui n'est pas garanti.
}

\approfondissement{
Historiquement, la relation $\ln'(x) = \dfrac{1}{x}$ correspond au problème de la \textbf{quadrature de l'hyperbole} : l'aire sous la courbe de $y=\dfrac{1}{x}$ entre $1$ et $x$ définit une fonction dont la dérivée vaut $\dfrac{1}{x}$ — c'est ainsi que certains mathématiciens (Grégoire de Saint-Vincent, au dix-septième siècle) ont approché la notion de logarithme par le calcul intégral, avant même que le lien avec l'exponentielle ne soit pleinement formalisé.
}
